% =====================================================================
%  Shakespeare
%  A Receiver Language for Protein Processes:
%  the Cut Primitive, the Biomolecular Leaf Algebra, and the
%  Causal Propagation Table as a Performable Play
%
%  Language specification. The framework results it binds to are
%  stated as named operations, not re-derived.
% =====================================================================
\documentclass[11pt,a4paper]{article}

\usepackage[utf8]{inputenc}
\usepackage[T1]{fontenc}
\usepackage{lmodern}
\usepackage[a4paper,margin=2.4cm]{geometry}
\usepackage{amsmath,amssymb,amsthm,mathtools}
\usepackage{bm}
\usepackage{mathrsfs}
\usepackage{booktabs}
\usepackage{array}
\usepackage{enumitem}
\usepackage{microtype}
\usepackage{xcolor}
\usepackage{listings}
\usepackage[round,sort&compress]{natbib}
\usepackage[colorlinks=true,linkcolor=blue!45!black,
            citecolor=blue!45!black,urlcolor=blue!45!black]{hyperref}
\usepackage{cleveref}

% ---- Theorem-like environments ----
\newtheoremstyle{thmstyle}{\topsep}{\topsep}{\itshape}{}{\bfseries}{.}{.5em}{}
\theoremstyle{thmstyle}
\newtheorem{theorem}{Theorem}[section]
\newtheorem{lemma}[theorem]{Lemma}
\newtheorem{proposition}[theorem]{Proposition}
\newtheorem{corollary}[theorem]{Corollary}
\theoremstyle{definition}
\newtheorem{definition}[theorem]{Definition}
\newtheorem{rulename}[theorem]{Rule}
\newtheorem{principle}[theorem]{Principle}
\newtheorem{example}[theorem]{Example}
\theoremstyle{remark}
\newtheorem{remark}[theorem]{Remark}
\newtheorem{notation}[theorem]{Notation}

% ---- shakespeare listing style ----
\definecolor{shkw}{RGB}{30,64,175}
\definecolor{shcom}{RGB}{22,101,52}
\definecolor{shnum}{RGB}{180,83,9}
\definecolor{shstr}{RGB}{159,18,57}
\definecolor{shbg}{RGB}{250,250,250}
\lstdefinelanguage{shakespeare}{
  morekeywords={receiver,floor,cut,contact,close,complex,fold,track,until,
    yield,when,observe,catalyze,mutate,complete,where,in,as,let,medium,
    converge,diverge,with,by,reps,representation,switch,import,module,
    export,assert,from,to,if,else,scene,perform,leaf,tree},
  morekeywords=[2]{residue,cofactor,electronic,solvent,substrate,
    Atom,Contact,Complex,Path,Trajectory,Signature,Structure,Scalar,Cut,
    S,nu,thick,delta,valence,shape,vacancy,floorval,count,coherence,
    spin,orbital,depth},
  sensitive=true,
  morecomment=[l]{--},
  morestring=[b]",
  keywordstyle=\color{shkw}\bfseries,
  keywordstyle=[2]\color{shnum},
  commentstyle=\color{shcom}\itshape,
  stringstyle=\color{shstr},
  basicstyle=\ttfamily\small,
  backgroundcolor=\color{shbg},
  frame=single,framesep=4pt,xleftmargin=6pt,xrightmargin=4pt,
  breaklines=true,showstringspaces=false,columns=fullflexible,tabsize=2
}
\lstset{language=shakespeare}

\crefname{rulename}{Rule}{Rules}
\Crefname{rulename}{Rule}{Rules}
\crefname{principle}{Principle}{Principles}
\Crefname{principle}{Principle}{Principles}

% ---- shorthand ----
\newcommand{\R}{\mathbb{R}}
\newcommand{\N}{\mathbb{N}}
\newcommand{\flo}{\beta}
\newcommand{\res}{\rho}
\newcommand{\cut}{\partial}
\newcommand{\medium}{\mathsf{m}}
\newcommand{\Recv}{\mathcal{R}_{\mathrm{bio}}}
\newcommand{\Sexp}{\xi}
\newcommand{\Sig}{\Sigma}
\newcommand{\eval}{\mathrm{eval}}
\newcommand{\evalto}{\Downarrow}
\newcommand{\dd}{\,\mathrm{d}}
\newcommand{\bnf}{\;::=\;}
\newcommand{\alt}{\;\mid\;}
\newcommand{\shk}[1]{\texttt{#1}}
\newcommand{\abs}[1]{\lvert#1\rvert}
\newcommand{\Osc}{\mathfrak{O}}
\newcommand{\Cat}{\mathfrak{C}}
\newcommand{\Part}{\mathfrak{P}}

\title{\textbf{Shakespeare:\\[3pt]
A Receiver Language for Protein Processes---\\
the Cut Primitive, the Biomolecular Leaf Algebra,\\
and the Causal Propagation Table as a Performable Play}}

\author{%
Kundai Farai Sachikonye\\
\small Technical University of Munich\\
\small AIMe Registry for Artificial Intelligence\\
\small \texttt{kundai.sachikonye@wzw.tum.de}}

\date{\today}

\begin{document}
\maketitle

% =====================================================================
\begin{abstract}
\noindent
We specify \emph{Shakespeare} (source extension \shk{.shk}), a
domain-specific language whose programs denote \emph{protein processes}
---folding, contact, complex formation, catalysis, mutation, and
partial-observation completion---and whose evaluation is not a
simulation of those processes but a \emph{measurement} of them. The
language has a single computational primitive, the \emph{cut}: the act
of individuating a part of a bounded whole at a strictly positive,
irreducible cost, the floor $\flo>0$. Its value model is the
\emph{biomolecular receiver} $\Recv$: every Shakespeare value is an
S-expression over a five-class oscillator leaf algebra (residue,
cofactor, electronic, solvent, substrate), and every evaluation is a run
of the receiver's typed morphism chain
$\eval = \mathrm{access}\circ\mathrm{fuse}\circ\mathrm{catalyze}^{*}\circ
\mathrm{observe}$. Generation, contact, closure, folding, and process
tracking are one operation---the cut---at increasing arity, exactly as
the underlying framework requires: a reaction step and a measurement step
are the same cut, so \shk{catalyze} and \shk{observe} are one verb read
two ways, and \shk{track\,\ldots\,until\,converge\,yield} \emph{is} the
causal propagation table evaluated over proteins. We give: (i)~the
semantic model---the contact graph, the receiver, the floor, the leaf
algebra, and accountability, on which every value carries the floor at
which it was cut; (ii)~a complete lexical and context-free grammar in
EBNF; (iii)~a static accountability type system with the five leaf
refinements and the process types $\mathsf{Fold}$, $\mathsf{Complex}$,
$\mathsf{Path}$, $\mathsf{Trajectory}$, in which a value claiming zero
residue is rejected at compile time; (iv)~a small-step operational
semantics in which evaluation is measurement (rendering \emph{is}
observation), the committed-cut count is a monotone intrinsic clock, and
a track's intermediate steps are free while only terminal
$S$-entropy convergence is judged; (v)~a standard library binding each
verb to a validated operation of the framework
(\shk{fold}$\to$the Kuramoto/coupling-spectrum folding chain;
\shk{$\sim$}$\to$boundary/phase-lock contact; \shk{close}$\to$closure to a
receiver tree; \shk{track}$\to$accountable propagation with convergence
admissibility, path opacity, and representation mobility;
\shk{catalyze}$\to$the catalytic cycle under the Miracle Principle;
\shk{mutate}$\to$re-cut of the coupling graph; \shk{complete}$\to$
trajectory completion from partial omics); and (vi)~a two-target
compilation model. Shakespeare compiles, from one typed front end and one
back-end-agnostic Cut-IR, to a Rust reference runtime in which cuts are
exact maximum-flow minimum cuts, and to a WebGL\,2 runtime---the online
Shakespeare tool---in which cuts are GPU render passes and the framebuffer
that results \emph{is} the measurement. We prove the two back ends are
observationally equivalent up to the floor. Worked plays fold crambin,
run the seven-state cytochrome CYP3A4 catalytic cycle (with the Miracle
Principle licensing the transient Compound~I), and score a point mutation.
The protein is a protocol; a Shakespeare program is its performance.
\end{abstract}

\noindent\textbf{Keywords:} domain-specific language; protein folding;
biomolecular receiver; cut primitive; contact graph; leaf algebra;
causal propagation table; accountability type system; operational
semantics; minimum cut; rendering-as-measurement; two-target compilation;
cytochrome P450; Rust; WebGL.

\tableofcontents

% =====================================================================
\section{Introduction}
\label{sec:intro}
% =====================================================================

\subsection{Proteins as protocols, plays as performances}

A protein is not a static object; it is a process---a chain of
individuations that folds a sequence, brings residues into contact,
binds a substrate, turns it over, and returns. The framework underlying
this language has shown that each such process is, in form, an
\emph{operation}: a procedure that folds a chain by frequency matching
\citep{sachikonye2025fold}, that decides and shapes a contact by boundary
minimisation \citep{sachikonye2025chem}, that tracks an item's
uncertainty through a process and returns an accountable account of what
happened to it \citep{sachikonye2025causal}, and that evaluates any
biomolecular S-expression---polypeptide, cofactor, electron, solvent,
substrate---under one receiver $\Recv$ \citep{sachikonye2025expralg}. Each
of these results is already a program with a proof of correctness
attached.

\emph{Shakespeare} makes that literal. It is a language whose statements
\emph{are} those operations, so that a protein-process paper becomes a
script and a derivation becomes a run. We name it for what it does: a
receiver evaluates an \emph{expression}---an S-expression, a scripted
sequence of parts and their individuations---and to run the program is to
\emph{perform} it, not to depict it. A Shakespeare program is a play; its
execution is a performance that measures the protein rather than
simulating it (\Cref{prin:measure}).

The design discipline is that we implement only what is already
established. Every verb binds to an operation of the framework that has an
explicit construction and has been validated numerically; the language
adds composition, typing, and execution, not new biology.

\subsection{One primitive: the cut}

The language has a single computational primitive, the \emph{cut}, and a
single reason for it. To individuate any part of a bounded whole---to make
it a distinct thing---costs a strictly positive, irreducible minimum, the
floor $\flo>0$, itself forced by the non-completability of the whole
against which a part is told apart \citep{sachikonye2025graphs,%
sachikonye2025causal}. Every Shakespeare operation is a cut:

\begin{itemize}[leftmargin=1.4em,itemsep=2pt]
\item A \emph{leaf} is a cut of arity one: individuating one oscillator
leaf---a residue, cofactor, electron, solvent molecule, or
substrate---from the medium (\Cref{sec:stdlib-leaf}).
\item A \emph{contact} is a cut of arity two: the single shared boundary
between two leaves, realised when phase-locking (equivalently,
boundary-sharing) lowers total thickness (\Cref{sec:stdlib-contact}).
\item A \emph{fold} or \emph{complex} is a closure of cuts: cuts continued
until every constituent reaches its partition-closure minimum, yielding a
receiver tree (\Cref{sec:stdlib-fold}).
\item A \emph{track} is a residue-chained sequence of cuts: the residue of
each cut is the cause of the next, and the chain is the propagation of an
item's uncertainty through a process, admissible iff it $S$-entropy
converges to the observed output (\Cref{sec:stdlib-track}).
\end{itemize}

Folding is not a different verb from tracking; it is a track whose process
is the Kuramoto synchronisation of the chain's residue oscillators.
Catalysis is not a different verb from measurement; in the underlying
theory a reaction step and a measurement step are the \emph{same} cut,
differing only in whether the residue's causal link is resolved
\citep{sachikonye2025causal}. Shakespeare therefore has one verb (the cut)
and one composition rule (residue chaining); all other surface forms
---\shk{fold}, \shk{$\sim$}, \shk{close}, \shk{track}, \shk{catalyze},
\shk{mutate}, \shk{observe}, \shk{complete}---are sugar over these.

\subsection{The value model: the biomolecular receiver}

Where a general-purpose language computes on numbers, Shakespeare computes
on the biomolecular receiver $\Recv$ \citep{sachikonye2025expralg}. A
value is an S-expression over the five-class leaf algebra, and evaluation
is a run of the receiver's typed morphism chain
\begin{equation}
\eval_{\Recv}
= \mathrm{access}\circ\mathrm{fuse}\circ\mathrm{catalyze}^{*}
  \circ\mathrm{observe}.
\label{eq:chain-intro}
\end{equation}
The five leaf classes---residue, cofactor, electronic, solvent,
substrate---let one language address a whole protein system: the
polypeptide chain, its prosthetic groups, the electrons that transfer
through it, the ordered active-site waters, and the bound substrate. This
is the specific respect in which Shakespeare is a \emph{protein}-process
language and not merely a cheminformatics one: its leaves are the parts a
protein process actually moves.

\subsection{Two targets, one play}

Shakespeare exists in two parallel implementations from a single
specification:

\begin{enumerate}[leftmargin=1.6em,itemsep=2pt]
\item A \textbf{Rust} compiler and native runtime for the reference and
production version. Here a cut is realised as an exact maximum-flow
minimum-cut computation on the contact graph, and folding delegates to the
validated \shk{levinthal-folding} integrator. This is the authoritative
version.
\item A \textbf{TypeScript}/WebGL\,2 compiler and runtime---the online
Shakespeare tool. Here a cut is realised as a GPU render pass over the
contact-graph state encoded in textures: the framebuffer that results
\emph{is} the measurement, in keeping with the rendering--measurement
identity \citep{sachikonye2025sbs}. Only the domain observation hook is
protein-specific; the rest of the shader pipeline is universal.
\end{enumerate}

Both share one front end (lexer, parser, accountability type checker) and
one typed core intermediate representation, \emph{Cut-IR}
(\Cref{sec:ir}); they differ only in the back end that realises a cut. We
prove (\Cref{thm:target-equiv}) that the two back ends agree on every
program up to the floor---the one resolution at which the theory already
forbids finer distinction---so the online tool and the reference
implementation never disagree about anything the language can express.

\subsection{What this document specifies}

\Cref{sec:model} fixes the semantic model: the contact graph, the
receiver, the floor, the leaf algebra, and accountability.
\Cref{sec:lexical} gives the lexical grammar; \Cref{sec:grammar} the
context-free grammar in EBNF. \Cref{sec:types} defines the accountability
type system and proves soundness. \Cref{sec:opsem} gives the small-step
operational semantics and proves cut monotonicity and
measurement-not-simulation. \Cref{sec:stdlib} specifies the standard
library, binding each verb to its framework operation. \Cref{sec:ir}
defines Cut-IR and the two-target model and proves target equivalence.
\Cref{sec:examples} gives worked plays. \Cref{sec:impl} states
implementation requirements. Throughout, framework operations are named and
their input/output contracts stated; their derivations are external to
this document.

% =====================================================================
\section{The semantic model}
\label{sec:model}
% =====================================================================

Shakespeare is defined over a single mathematical object, the contact
graph carrying a receiver, and a single operation on it, the cut. We fix
these before any syntax.

\subsection{The contact graph and the medium}

\begin{definition}[Contact graph]
\label{def:contact-graph}
A \emph{contact graph} is a finite weighted graph $\Gamma=(V,E,w)$ with a
distinguished \emph{medium} vertex $\medium\in V$ adjacent to every other
vertex, and $w\colon E\to\R_{>0}$ a strictly positive weight. Vertices
$v\neq\medium$ are \emph{leaves}; an edge between two leaves is a
\emph{contact}; $w(e)$ is the cost of the separation it maintains. The
graph is the runtime state of a Shakespeare program: every value the
program holds is a structure on $\Gamma$.
\end{definition}

\begin{definition}[Floor]
\label{def:floor}
The \emph{floor} of a contact graph is a constant $\flo>0$ with
$w(e)\ge\flo$ for every edge. No separation is free. The floor is a
program-level parameter (set by \shk{floor}, \Cref{sec:grammar}); its
positivity is invariant and enforced by the type system
(\Cref{sec:types}). By the derivation of the framework
\citep{sachikonye2025graphs,sachikonye2025causal} the floor is not a free
posit but the irreducible residue a non-completable whole forces on each
finite part; for the biomolecular receiver it has the explicit
three-term estimate of \citep{sachikonye2025expralg}
(discretisation, oscillator $Q$-factor, conversion error).
\end{definition}

\begin{definition}[Cut]
\label{def:cut}
A \emph{cut} of a leaf set $S$ ($\varnothing\neq S\subsetneq V$,
$\medium\notin S$) is the edge boundary
$\cut(S)=\{e\in E:\abs{e\cap S}=1\}$, of weight
$w(\cut(S))=\sum_{e\in\cut(S)}w(e)\ge\flo$. The cut individuates $S$ from
the rest. A cut \emph{event} commits one such boundary: it fixes the edges
of $\cut(S)$ as realised separations and advances the program's cut count
by one.
\end{definition}

\begin{definition}[Residue]
\label{def:residue}
The \emph{residue} of a cut event is the weight $\res=w(\cut(S))\ge\flo$ it
deposits as committed boundary. The residue is the value the cut produces
and the cause of any subsequent cut: a chain of cuts is one in which each
residue is consumed as the boundary condition of the next
(\Cref{sec:stdlib-track}).
\end{definition}

\subsection{The biomolecular receiver and its leaves}

\begin{definition}[Receiver]
\label{def:receiver}
A \emph{receiver} is a tuple
$\Recv=(\mathcal{H}, \eval_{\Recv}, \{\Phi^{R\to R'}\}, \flo(\Recv))$ where
$\mathcal{H}$ is the representation space, $\eval_{\Recv}$ the evaluation
map (\Cref{eq:chain-intro}), $\Phi^{R\to R'}$ the conversion functors
between the oscillatory, categorical, and partition representations
$\{\Osc,\Cat,\Part\}$, and $\flo(\Recv)>0$ the intrinsic floor. The
functors form a closed cycle under $S_3$ action, so the receiver is
representation-coherent: $\eval_{\Recv}$ is invariant under cyclic
conversion \citep{sachikonye2025expralg}. Shakespeare fixes the receiver
by the declaration \shk{receiver\ bio}.
\end{definition}

\begin{definition}[Oscillator leaf; five classes]
\label{def:leaf}
An \emph{oscillator leaf} is a tuple
$\ell=(\omega,\mathbf{S},\phi,R,\kappa)$ with natural frequency
$\omega\in\R_{>0}$, S-entropy coordinate
$\mathbf{S}=(S_k,S_t,S_e)\in[0,1]^3$, instantaneous phase
$\phi\in[0,2\pi)$, representation tag $R\in\{\Osc,\Cat,\Part\}$, and class
\[
\kappa\in\{\mathrm{res},\mathrm{cof},\mathrm{ele},\mathrm{sol},\mathrm{sub}\}:
\]
\emph{residue} (an amino acid, $\omega$ from the amide-I band shifted by
side-chain identity), \emph{cofactor} (a prosthetic group: heme, FAD, FMN,
NADPH, bound $\mathrm{O_2}$), \emph{electronic} (an electron at partition
coordinate $(n,\ell,m,s)$), \emph{solvent} (an ordered active-site water),
and \emph{substrate} (a bound small-molecule ligand). Each leaf is a cut
of arity one; its S-entropy coordinate is its address and its residue is
$w(\cut(\{\ell\}))\ge\flo$.
\end{definition}

\begin{principle}[Accountability]
\label{prin:accountability}
Every Shakespeare value is a cut, and every cut carries its residue. There
are no bare numbers in Shakespeare: a quantity is always paired with the
floor at which it was individuated. A value claiming residue below the
floor, or zero residue, is not a value---it is the forbidden sharp cut, and
the type system rejects it (\Cref{thm:no-zero-value}).
\end{principle}

\begin{remark}[Why a receiver and not a number line]
\label{rem:why-receiver}
Shakespeare computes on the receiver rather than on real numbers because
the quantities it manipulates---a residue's fold state, a contact's
existence, a catalytic path's accountability---are minimum cuts and
partition signatures, not scalars. A fold is the receiver tree that closes
a chain's vacancies; a contact is the invariant minimum cut between two
leaves; a track's result is an accountable chain of cuts. Representing
these as receiver structures (rather than as fitted scalars) is what lets
the single cut primitive serve leaf generation, contact, folding, and
tracking alike, and is what the two back ends each realise: exactly (Rust,
max-flow) or by rendering (WebGL, GPU).
\end{remark}

\subsection{Measurement, not simulation}

\begin{principle}[Evaluation is measurement]
\label{prin:measure}
Executing a Shakespeare cut performs the individuation; it does not
simulate a description of it. The result of \shk{fold\ "1CRN"} is the
measured partition signature of crambin's fold, carried as an accountable
value, not a numerical trajectory toward it. Operationally
(\Cref{sec:opsem}) a cut evaluation mutates the contact graph---it commits
boundary and advances the count---rather than returning a pure function of
inputs. The distinction is observable: re-evaluating a cut is a new cut at
a higher count, never a cached recomputation (\Cref{thm:monotone}). On the
WebGL target this is literal: the render pass that produces the framebuffer
\emph{is} the observation \citep{sachikonye2025sbs}.
\end{principle}

% =====================================================================
\section{Lexical structure}
\label{sec:lexical}
% =====================================================================

A Shakespeare source file is UTF-8 text with extension \shk{.shk}. Lexing
produces a token stream from the following classes.

\begin{definition}[Token classes]
\label{def:tokens}
\leavevmode
\begin{itemize}[leftmargin=1.4em,itemsep=1pt]
\item \textbf{Comments.} \shk{--} to end of line; discarded.
\item \textbf{Keywords.} \shk{receiver}, \shk{floor}, \shk{cut},
\shk{contact}, \shk{close}, \shk{complex}, \shk{fold}, \shk{track},
\shk{until}, \shk{yield}, \shk{when}, \shk{observe}, \shk{catalyze},
\shk{mutate}, \shk{complete}, \shk{in}, \shk{as}, \shk{let},
\shk{medium}, \shk{converge}, \shk{diverge}, \shk{with}, \shk{by},
\shk{reps}, \shk{import}, \shk{module}, \shk{export}, \shk{assert},
\shk{scene}, \shk{perform}.
\item \textbf{Leaf-class words.} \shk{residue}, \shk{cofactor},
\shk{electronic}, \shk{solvent}, \shk{substrate}.
\item \textbf{Identifiers.} \shk{[A-Za-z\_][A-Za-z0-9\_]*}, not a keyword.
\item \textbf{Numbers.} decimal integers and floats with optional
exponent: \shk{[0-9]+(.[0-9]+)?([eE][+-]?[0-9]+)?}.
\item \textbf{Strings.} double-quoted, with the usual escapes; used for
PDB identifiers, sequences, and file references.
\item \textbf{Operators and punctuation.} \shk{:=} (bind), \shk{\~{}} (cut
between / contact), \shk{( ) [ ] \{ \} , : . > < >= <= == ->}.
\end{itemize}
Whitespace is insignificant except as a token separator; Shakespeare is
not layout-sensitive.
\end{definition}

\begin{notation}[Floors in literals]
\label{not:literals}
A numeric literal may carry an explicit floor annotation with the suffix
form \shk{value\#floor} (read ``value at floor''); e.g.\ \shk{0.8\#1e-4}.
A literal without an explicit floor inherits the ambient program floor set
by the most recent \shk{floor} declaration in scope. There is no way to
write a literal of zero floor (\Cref{thm:no-zero-value}).
\end{notation}

% =====================================================================
\section{Grammar}
\label{sec:grammar}
% =====================================================================

We give the context-free grammar in EBNF \citep{ebnf2017}. Nonterminals
are \textit{italicised}; terminals are \shk{typewriter}; \shk{\{x\}} is
zero or more, \shk{[x]} optional, \shk{|} alternation.

\begin{definition}[Shakespeare grammar]
\label{def:grammar}
\begin{align*}
\textit{program} &\bnf \{\textit{decl}\} \\[2pt]
\textit{decl} &\bnf \textit{receiverDecl} \alt \textit{floorDecl}
   \alt \textit{import} \alt \textit{moduleDecl} \alt \textit{scene}
   \alt \textit{stmt} \\[2pt]
\textit{receiverDecl} &\bnf \shk{receiver}\ \textit{ident} \\
\textit{floorDecl} &\bnf \shk{floor}\ \textit{number} \\
\textit{import} &\bnf \shk{import}\ \textit{qname} \\
\textit{moduleDecl} &\bnf \shk{module}\ \textit{ident}\ \shk{\{}\
   \{\textit{decl}\}\ \shk{\}} \\
\textit{scene} &\bnf \shk{scene}\ \textit{ident}\ \shk{\{}\
   \{\textit{stmt}\}\ \shk{\}} \\[2pt]
\textit{stmt} &\bnf \textit{bind} \alt \textit{track} \alt \textit{observe}
   \alt \textit{assert} \alt \textit{perform} \alt \textit{expr} \\[2pt]
\textit{bind} &\bnf [\shk{let}]\ \textit{ident}\ \shk{:=}\ \textit{expr} \\
\textit{perform} &\bnf \shk{perform}\ \textit{ident} \\[2pt]
\textit{expr} &\bnf \textit{leafExpr} \alt \textit{contactExpr}
   \alt \textit{closeExpr} \alt \textit{foldExpr}
   \alt \textit{catalyzeExpr} \alt \textit{mutateExpr}
   \alt \textit{completeExpr} \\
   &\qquad \alt \textit{call} \alt \textit{ident} \alt \textit{number}
   \alt \textit{string} \alt \shk{(}\ \textit{expr}\ \shk{)} \\[2pt]
\textit{leafExpr} &\bnf \shk{cut}\ \textit{cutArg}
   \alt \textit{leafClass}\ \textit{leafArg} \\
\textit{leafClass} &\bnf \shk{residue} \alt \shk{cofactor}
   \alt \shk{electronic} \alt \shk{solvent} \alt \shk{substrate} \\
\textit{cutArg} &\bnf \textit{number} \alt \textit{string}
   \alt \textit{ident} \\
\textit{leafArg} &\bnf \textit{string} \alt \textit{ident}
   \alt \textit{number} \\[2pt]
\textit{contactExpr} &\bnf \textit{expr}\ \shk{\~{}}\ \textit{expr}\
   [\,\shk{when}\ \textit{cond}\,] \\[2pt]
\textit{closeExpr} &\bnf (\shk{close}\alt\shk{complex})\ \textit{ident}\
   \shk{(}\ \textit{argList}\ \shk{)}\ [\,\shk{by}\ \textit{ident}\,] \\[2pt]
\textit{foldExpr} &\bnf \shk{fold}\ (\textit{string}\alt\textit{ident})\
   [\,\shk{with}\ \textit{argList}\,] \\[2pt]
\textit{catalyzeExpr} &\bnf \shk{catalyze}\ \textit{ident}\ \shk{in}\
   \textit{expr}\ [\,\shk{yield}\ \textit{ident}\,] \\[2pt]
\textit{mutateExpr} &\bnf \shk{mutate}\ \textit{expr}\ \shk{by}\
   \textit{mutSpec} \\
\textit{mutSpec} &\bnf \textit{ident}\ \shk{->}\ \textit{ident}\
   [\,\shk{at}\ \textit{number}\,] \\[2pt]
\textit{completeExpr} &\bnf \shk{complete}\ \textit{expr}\ \shk{from}\
   \textit{argList} \\[2pt]
\textit{track} &\bnf \shk{track}\ \textit{ident}\ \shk{in}\ \textit{expr}
   \\
   &\qquad [\,\shk{with}\ \shk{reps}\ \textit{repList}\,]\ \\
   &\qquad \shk{until}\ \textit{admit}\ \ \shk{yield}\ \textit{ident} \\[2pt]
\textit{admit} &\bnf \shk{converge} \alt \shk{diverge} \alt \textit{cond}
   \\[2pt]
\textit{observe} &\bnf \shk{observe}\ \textit{expr}\
   [\,\shk{as}\ \textit{ident}\,] \\[2pt]
\textit{assert} &\bnf \shk{assert}\ \textit{cond}\
   [\,\shk{emit}\ \textit{string}\,] \\[2pt]
\textit{call} &\bnf \textit{qname}\ \shk{(}\ [\textit{argList}]\ \shk{)} \\
\textit{argList} &\bnf \textit{arg}\ \{\shk{,}\ \textit{arg}\} \\
\textit{arg} &\bnf [\textit{ident}\ \shk{:}]\ \textit{expr} \\[2pt]
\textit{cond} &\bnf \textit{expr}\ \textit{relop}\ \textit{expr} \\
\textit{relop} &\bnf \shk{>} \alt \shk{<} \alt \shk{>=} \alt \shk{<=}
   \alt \shk{==} \\[2pt]
\textit{repList} &\bnf \textit{ident}\ \{\shk{,}\ \textit{ident}\} \\
\textit{qname} &\bnf \textit{ident}\ \{\shk{.}\ \textit{ident}\} \\
\textit{number} &\bnf \textit{numlit}\ [\shk{\#}\ \textit{numlit}]
\end{align*}
\end{definition}

\begin{remark}[Reading the core forms]
\label{rem:core-forms}
The whole language is small. \shk{cut\ N} or \shk{residue\ "ALA"}
individuates a leaf (arity-one cut); \shk{a\ \~{}\ b\ when\ c} contacts two
leaves (arity-two cut, guarded); \shk{close\ X(\ldots)} /
\shk{complex\ X(\ldots)} drives constituents to closure;
\shk{fold\ "1CRN"} runs the folding chain; \shk{track\ x\ in\ P\ until\
converge\ yield\ r} propagates and binds the accountable result;
\shk{catalyze\ s\ in\ E} turns a substrate over in an enzyme;
\shk{mutate\ P\ by\ A->B\ at\ i} re-cuts the coupling graph;
\shk{complete\ P\ from\ obs} fills the state from partial observation.
\shk{observe} forces measurement now; \shk{assert} is a pre-flight check;
\shk{scene} groups statements into a named act; \shk{perform} runs one.
Everything else---\shk{let}, \shk{import}, \shk{module}, named
arguments---is conventional.
\end{remark}

% =====================================================================
\section{The accountability type system}
\label{sec:types}
% =====================================================================

Shakespeare's types record \emph{what kind of cut} a value is and \emph{at
what floor}. The system's purpose is to make
Principle~\ref{prin:accountability} statically enforceable: a value of
zero residue cannot be typed.

\subsection{Types}

\begin{definition}[Types]
\label{def:types}
The Shakespeare types are
\[
\tau \bnf \mathsf{Leaf}_\kappa \alt \mathsf{Contact} \alt \mathsf{Complex}
   \alt \mathsf{Fold} \alt \mathsf{Path} \alt \mathsf{Trajectory}
   \alt \mathsf{Signature} \alt \mathsf{Scalar} \alt \mathsf{Cut},
\]
where $\kappa\in\{\mathrm{res},\mathrm{cof},\mathrm{ele},\mathrm{sol},
\mathrm{sub}\}$. $\mathsf{Cut}$ is the supertype of all individuated
values. $\mathsf{Leaf}_\kappa$ is an arity-one cut of class $\kappa$;
$\mathsf{Contact}$ an arity-two cut; $\mathsf{Complex}$ and $\mathsf{Fold}$
are closures of cuts (a bound assembly and a folded chain respectively);
$\mathsf{Path}$ is a residue chain (the amalgamation of a track);
$\mathsf{Trajectory}$ is a time-indexed sequence of partition cells (a
catalytic or electron-transfer path); $\mathsf{Signature}$ is a partition
signature ($N\times N$ grid of $(n,\ell,m,s)$). $\mathsf{Scalar}$ is a
measured number; it is never bare. Every typed value $v$ additionally
carries a \emph{floor annotation} $\flo(v)>0$ and a \emph{residue}
$\res(v)\ge\flo(v)$.
\end{definition}

\begin{definition}[Typing judgement]
\label{def:judgement}
A judgement $\Gamma\vdash e:\tau\,@\,\flo$ reads ``in floor context
$\Gamma$ (which fixes the ambient floor and receiver), expression $e$ has
type $\tau$ and is individuated at floor $\flo>0$.'' The context carries
the ambient floor set by the most recent \shk{floor} declaration, the
active receiver set by \shk{receiver}, and the types of bound identifiers.
\end{definition}

\subsection{The positivity rule}

\begin{rulename}[Floor positivity]
\label{rule:positive}
Every value-introducing rule has the side condition $\flo>0$ and
$\res\ge\flo$. A literal or expression whose annotation forces $\flo\le0$
or $\res<\flo$ is ill-typed.
\[
\frac{\Gamma\vdash\textit{numlit}\;n \quad \flo>0}
     {\Gamma\vdash n\#\flo:\mathsf{Scalar}\,@\,\flo}
\qquad
\frac{\kappa\in\{\mathrm{res},\dots,\mathrm{sub}\}}
     {\Gamma\vdash \kappa\ a:\mathsf{Leaf}_\kappa\,@\,\flo_\Gamma}
\]
where $\flo_\Gamma$ is the ambient floor.
\end{rulename}

\begin{theorem}[No zero-residue value]
\label{thm:no-zero-value}
There is no well-typed Shakespeare expression of residue $0$. Equivalently,
the sharp cut is not expressible: every typeable value has
$\res\ge\flo>0$.
\end{theorem}

\begin{proof}
By induction on typing derivations. The only value introductions are
literals (\Cref{rule:positive}, side condition $\flo>0$, and a literal's
residue is its annotated floor), leaves (residue
$=w(\cut(\{\ell\}))\ge\flo>0$ by \Cref{def:cut,def:leaf}), and the derived
forms (contact, close, fold, track, catalyze, mutate, complete) whose
residues are sums or chains of cut residues, hence $\ge\flo$. No rule
introduces a value of residue $<\flo$, and $\flo>0$ is a side condition on
every floor-introduction. Hence no derivation concludes $\res=0$.
\end{proof}

\subsection{Verb typing}

\begin{rulename}[Leaf; contact; close/complex; fold]
\label{rule:verbs}
\[
\frac{\Gamma\vdash a:\mathsf{Leaf}_{\kappa_a}\quad
      \Gamma\vdash b:\mathsf{Leaf}_{\kappa_b}}
     {\Gamma\vdash a\ \shk{\~{}}\ b:\mathsf{Contact}\,@\,\flo_\Gamma}
\qquad
\frac{\Gamma\vdash X:\mathsf{Leaf}_\kappa\quad
      \Gamma\vdash a_i:\mathsf{Leaf}_{\kappa_i}}
     {\Gamma\vdash \shk{close}\ X(a_1,\dots,a_k)
      :\mathsf{Complex}\,@\,\flo_\Gamma}
\]
\[
\frac{\Gamma\vdash s:\mathsf{Scalar}\ \text{or}\ \mathsf{Leaf}_\kappa\
      \ (\text{PDB id, sequence, or chain})}
     {\Gamma\vdash \shk{fold}\ s:\mathsf{Fold}\,@\,\flo_\Gamma}
\]
The contact rule additionally requires its guard, when present, to be a
$\mathsf{Bool}$-valued \textit{cond}; an unguarded contact is admitted only
if static analysis cannot exclude a positive shared content
$\Delta\mathrm{thick}>0$ (\Cref{sec:stdlib-contact}).
\end{rulename}

\begin{rulename}[Track; catalyze; mutate; complete]
\label{rule:track}
\[
\frac{\Gamma\vdash x:\mathsf{Leaf}_\kappa\ \text{or}\ \mathsf{Complex}\quad
      \Gamma\vdash P:\mathsf{Complex}\ \text{or}\ \mathsf{Fold}\
        \text{or}\ \mathsf{Path}}
     {\Gamma\vdash \shk{track}\ x\ \shk{in}\ P\ \shk{until}\
        \textit{admit}\ \shk{yield}\ r:\mathsf{Path}\,@\,\flo_\Gamma}
\]
\[
\frac{\Gamma\vdash s:\mathsf{Leaf}_{\mathrm{sub}}\quad
      \Gamma\vdash E:\mathsf{Complex}\ \text{or}\ \mathsf{Fold}}
     {\Gamma\vdash \shk{catalyze}\ s\ \shk{in}\ E
      :\mathsf{Trajectory}\,@\,\flo_\Gamma}
\]
\[
\frac{\Gamma\vdash P:\mathsf{Fold}\ \text{or}\ \mathsf{Complex}}
     {\Gamma\vdash \shk{mutate}\ P\ \shk{by}\ m:\mathsf{Fold}\,@\,\flo_\Gamma}
\qquad
\frac{\Gamma\vdash P:\mathsf{Fold}\quad \Gamma\vdash o_i:\mathsf{Signature}}
     {\Gamma\vdash \shk{complete}\ P\ \shk{from}\ \bar o
      :\mathsf{Signature}\,@\,\flo_\Gamma}
\]
The tracked result $r$ has type $\mathsf{Path}$; its residue is the total
committed boundary of the propagation (\Cref{sec:stdlib-track}). Catalysis
yields a $\mathsf{Trajectory}$---the sequence of partition cells the
substrate traverses---read either as reaction or as measurement
(\Cref{sec:stdlib-catalyze}).
\end{rulename}

\begin{theorem}[Type soundness]
\label{thm:soundness}
Shakespeare typing is sound for the operational semantics of
\Cref{sec:opsem}: if $\Gamma\vdash e:\tau\,@\,\flo$ and $e$ evaluates to a
value $v$, then $v$ has type $\tau$, floor $\flo>0$, and residue
$\res(v)\ge\flo$ (progress and preservation).
\end{theorem}

\begin{proof}[Proof sketch]
Standard syntactic soundness \citep{pierce2002}. \emph{Progress}: a
well-typed closed expression is a value or steps (every form has a
reduction rule in \Cref{sec:opsem}). \emph{Preservation}: each reduction
rule preserves the type and floor annotation, and---by
\Cref{thm:no-zero-value} applied at each step---preserves $\res\ge\flo>0$.
The only non-standard clause is that reductions are state-mutating
(measurement, \Cref{prin:measure}); preservation holds because the cut
count and committed boundary are monotone (\Cref{thm:monotone}), so no step
lowers a residue below the floor. Type safety across the five leaf classes
follows because every operation after \shk{observe} acts on the uniform
partition signature, into which class information is folded during the
observe step \citep{sachikonye2025expralg}.
\end{proof}

% =====================================================================
\section{Operational semantics}
\label{sec:opsem}
% =====================================================================

We give a small-step semantics over configurations
$\langle\Gamma_G,M,e\rangle$, where $\Gamma_G$ is the current contact graph
(\Cref{def:contact-graph}), $M\in\N$ is the committed-cut count (the
program clock), and $e$ is the expression under evaluation. We write
$\langle\Gamma_G,M,e\rangle\to\langle\Gamma_G',M',v\rangle$ for one step.
Evaluation is measurement: steps mutate $\Gamma_G$ and increment $M$.

\subsection{Leaf and morphism-chain reduction}

\begin{rulename}[Leaf commits and counts]
\label{rule:leaf-step}
\[
\frac{S=\textsc{Individuate}(\Gamma_G,\kappa,a)\qquad
      \res=w(\cut(S))\ge\flo}
     {\langle\Gamma_G,M,\kappa\ a\rangle \to
      \langle\Gamma_G\!\oplus\!\cut(S),\,M+1,\,
      v_{\mathsf{Leaf}_\kappa}(S,\res)\rangle}
\]
where \textsc{Individuate} is the framework's leaf-individuation operation
(\Cref{sec:stdlib-leaf}), $\Gamma_G\!\oplus\!\cut(S)$ is $\Gamma_G$ with the
cut's boundary committed, and $v_{\mathsf{Leaf}_\kappa}(S,\res)$ is the
resulting accountable leaf value.
\end{rulename}

\begin{rulename}[Closure runs the morphism chain]
\label{rule:chain-step}
\[
\frac{\Sig=(\mathrm{access}\circ\mathrm{fuse}\circ\mathrm{catalyze}^{*}
      \circ\mathrm{observe})(\textsc{Tree}(e))\qquad
      \res=\textstyle\sum_i w(\cut(S_i))\ge\flo}
     {\langle\Gamma_G,M,\shk{close}\ e\rangle \to
      \langle\Gamma_G',\,M+k,\,v_{\mathsf{Complex}}(\Sig,\res)\rangle}
\]
where $\textsc{Tree}(e)$ is the receiver tree of $e$, the chain is
evaluated in the sense of \Cref{eq:chain-intro}, and $k$ is the number of
cut events it commits (the closure count). The \shk{fold} rule is
identical with $\textsc{Tree}$ built from the Kuramoto synchronisation of
the chain's residue leaves (\Cref{sec:stdlib-fold}).
\end{rulename}

The contact, catalyze, mutate, and complete rules are analogous: each
commits boundary and increments $M$ by the number of cut events it
performs (one for a contact, the closure count for \shk{close}/\shk{fold},
the chain length for \shk{track}/\shk{catalyze}).

\begin{theorem}[Cut monotonicity / intrinsic clock]
\label{thm:monotone}
For any program, the committed-cut count $M$ is non-decreasing along
evaluation and strictly increases at every cut event. No reduction
decreases $M$. Consequently re-evaluating the same surface expression is a
new cut at a strictly higher count, never a cached recomputation, and two
configurations with identical graphs but different $M$ are distinct
states.
\end{theorem}

\begin{proof}
Every value-producing rule (\Cref{rule:leaf-step,rule:chain-step} and their
analogues) increments $M$; no rule decrements it (there is no reduction
that un-commits boundary---``undo'' is itself a further cut). Hence $M$ is
monotone and strictly increases per cut event. Distinctness of equal-graph
states at different $M$ is immediate since $M$ is part of the
configuration.
\end{proof}

\begin{corollary}[Measurement, not simulation]
\label{cor:measure}
Because evaluation mutates $\Gamma_G$ and advances $M$, a Shakespeare cut
cannot be a pure function memoised across calls: each occurrence performs
the individuation afresh and is recorded in the clock. This realises
Principle~\ref{prin:measure}: performing the play \emph{is} performing the
measurements it names.
\end{corollary}

\subsection{Track reduction and convergence}

\begin{rulename}[Track is a residue chain]
\label{rule:track-step}
\[
\frac{\begin{array}{c}
       (S_1,\dots,S_k)=\textsc{Propagate}(\Gamma_G,x,P)\quad
       \res_i=w(\cut(S_i))\\
       \textsc{Admit}(\res_1,\dots,\res_k)=\textit{admit}
      \end{array}}
     {\langle\Gamma_G,M,\shk{track}\ x\ \shk{in}\ P\ \shk{until}\
        \textit{admit}\ \shk{yield}\ r\rangle
      \to \langle\Gamma_G',\,M+k,\,
      v_{\mathsf{Path}}(S_{1:k},\textstyle\sum_i\res_i)\rangle}
\]
\textsc{Propagate} is the framework's accountable-propagation operation
\citep{sachikonye2025causal}; each $\res_i$ is the residue of cut $i$ and
the boundary condition of cut $i{+}1$ (residue chaining). \textsc{Admit}
tests the convergence criterion ($S$-entropy convergence for
\shk{converge}); a non-admissible chain reduces to the empty path (no
yield).
\end{rulename}

\begin{remark}[Local steps are free; only convergence is judged]
\label{rem:local-free}
In keeping with the underlying propagation theory
\citep{sachikonye2025causal}, the intermediate cuts of a track are
unconstrained: a step may be locally extreme (even ``impossible'')
provided the chain composes to a value that meets the admissibility
criterion. \textsc{Admit} therefore inspects only the terminal convergence
of the chain, not the plausibility of intermediate $\res_i$. This is what
makes \shk{track\,\ldots\,until\,converge} a search over admissible
propagations rather than a single forced path, and it is what licenses a
catalytic track through a transient intermediate such as Compound~I
(\Cref{sec:stdlib-catalyze}, the Miracle Principle). Optional
\shk{with\ reps\ }$R$ permits representation switching between steps
(mean-recovery-equivalent decompositions) at no additional cut cost
(representation mobility, \citealt{sachikonye2025causal}).
\end{remark}

% =====================================================================
\section{The standard library}
\label{sec:stdlib}
% =====================================================================

Each Shakespeare verb binds to a named operation of the underlying
framework. We state the contract (inputs, output, the operation it
invokes); the operation's construction and validation are external to this
document.

\subsection{\texorpdfstring{Leaves: \shk{cut}, \shk{residue}, \shk{cofactor}, \ldots}{Leaves}}
\label{sec:stdlib-leaf}

\begin{definition}[\textsc{Individuate}]
\label{def:individuate}
A leaf form invokes \textsc{Individuate}$(\kappa,a)$: from a class $\kappa$
and an identity $a$, assign the leaf its natural frequency $\omega$,
S-entropy coordinate $\mathbf{S}\in[0,1]^3$, and---for electronic
leaves---partition coordinate $(n,\ell,m,s)$ under the shell capacity
$C(n)=2n^2$, returning a $\mathsf{Leaf}_\kappa$ value carrying its address,
vacancy, and residue $\res=w(\cut(\{\ell\}))\ge\flo$
\citep{sachikonye2025expralg,sachikonye2025chem}. Residue leaves derive
$\omega$ from the amide-I band; cofactor, solvent, and substrate leaves
from the dominant vibrational mode of the group; electronic leaves from the
atomic transition frequency. E.g.\ \shk{residue\ "TRP"} yields the
tryptophan leaf; \shk{cofactor\ "heme"} the heme sub-tree.
\end{definition}

\subsection{\texorpdfstring{\shk{$\sim$} --- contact}{Contact}}
\label{sec:stdlib-contact}

\begin{definition}[\textsc{Contact}]
\label{def:contact}
\shk{a\ \~{}\ b} invokes \textsc{Contact}$(a,b)$: form the shared-boundary
cut between leaves $a,b$ and return a $\mathsf{Contact}$ value iff the
shared content $\Delta\mathrm{thick}(a,b)>0$---equivalently, iff their
oscillators phase-lock (angular level $\ell=0$) with coupling above
threshold \citep{sachikonye2025fold,sachikonye2025chem}. The optional guard
\shk{when\ }$c$ requires $c$ to hold (e.g.\ \shk{when\ delta\ >\ 0}). A
closed-shell or non-locking pair contacts with nothing.
\end{definition}

\subsection{\texorpdfstring{\shk{close}, \shk{complex}, \shk{fold} --- closure}{Closure}}
\label{sec:stdlib-fold}

\begin{definition}[\textsc{Close}, \textsc{Fold}]
\label{def:fold}
\shk{close\ X(a$_1$,\ldots,a$_k$)} / \shk{complex\ X(\ldots)} invokes
\textsc{Close}: form the set of contacts that drives every constituent's
vacancy to zero, returning a $\mathsf{Complex}$ value carrying the
stoichiometry, per-centre geometry, and total residue
\citep{sachikonye2025chem}. \shk{fold\ s} invokes \textsc{Fold}: run the
folding morphism chain on the chain $s$---Kuramoto synchronisation of the
residue-oscillator leaves, DFT of the coupling time series to a coupling
spectrum, then
$\mathrm{access}\circ\mathrm{fuse}\circ\mathrm{catalyze}^{*}\circ
\mathrm{observe}$---and return a $\mathsf{Fold}$ value carrying the
predicted contact map with secondary-structure annotation
\citep{sachikonye2025fold}. The native state is the unique global minimum
of the Kuramoto energy, so \shk{fold} halts at a floor-limited fixed point,
not a $3^N$ search.
\end{definition}

\subsection{\texorpdfstring{\shk{track} --- accountable propagation}{Track}}
\label{sec:stdlib-track}

\begin{definition}[\textsc{Propagate}, \textsc{Admit}]
\label{def:propagate}
\shk{track\ x\ in\ P\ until\ converge\ yield\ r} invokes
\textsc{Propagate}$(x,P)$: track item $x$'s individuation through process
$P$ as a residue-chained sequence of cuts, and \textsc{Admit} the chain iff
it $S$-entropy-converges to the output (the only admissibility test; local
steps free, \Cref{rem:local-free}). The yielded $\mathsf{Path}$ is the
\emph{amalgamation}: the accountable set of cuts---reaction steps and
measurement steps indifferently---by which $x$ reached the output
\citep{sachikonye2025causal}. This is the causal propagation table, now
over proteins; the resting catalogue (the folded fixed points) is
\shk{track} evaluated at rest.
\end{definition}

\subsection{\texorpdfstring{\shk{catalyze} --- reaction as measurement}{Catalyze}}
\label{sec:stdlib-catalyze}

\begin{definition}[\textsc{Catalyze}]
\label{def:catalyze}
\shk{catalyze\ s\ in\ E\ yield\ r} invokes \textsc{Catalyze}: track
substrate leaf $s$ through the enzyme complex $E$ to product as a
$\mathsf{Trajectory}$---a sequence of partition cells---read as a reaction.
Because a reaction step and a measurement step are the same cut
\citep{sachikonye2025causal}, \shk{catalyze} and \shk{observe} are one verb
read two ways: \shk{catalyze} leaves the residue's causal link unresolved
(entropy reading), \shk{observe} resolves it (information reading). The
Miracle Principle \citep{sachikonye2025expralg} licenses transient
intermediates: a catalytic sub-step may be locally maximally infeasible
(e.g.\ cytochrome-P450 Compound~I) provided the whole cycle closes
(substrate$\to$product in one enzyme conformation), so \shk{catalyze}
admits the full seven-state cycle without violating the framework's
selection rules.
\end{definition}

\subsection{\texorpdfstring{\shk{mutate} --- re-cut of the coupling graph}{Mutate}}
\label{sec:stdlib-mutate}

\begin{definition}[\textsc{Mutate}]
\label{def:mutate}
\shk{mutate\ P\ by\ A->B\ at\ i} invokes \textsc{Mutate}: replace the
residue leaf at position $i$, re-cut the coupling matrix $K_{ij}$, and
return the mutant $\mathsf{Fold}$. The variant effect is read from the
change in the coupling eigenstructure and in the folding convergence
\citep{sachikonye2025fold}: a destabilising mutation is one whose re-cut
graph fails to reach the closed fixed point or shifts the eigenvalue count
of secondary-structure blocks. This is an $O(N^2)$ query against the folded
receiver, not a re-fold from scratch.
\end{definition}

\subsection{\texorpdfstring{\shk{complete} --- trajectory completion}{Complete}}
\label{sec:stdlib-complete}

\begin{definition}[\textsc{Complete}]
\label{def:complete}
\shk{complete\ P\ from\ o$_1$,\ldots,o$_m$} invokes \textsc{Complete}:
given partial observations (omics signatures $o_i$ on a subset of leaves),
propagate the receiver's constraints---the circuit's Kirchhoff balance,
loop consistency, and backward trajectories---to a self-consistent complete
$\mathsf{Signature}$ over all leaves \citep{sachikonye2025sbs}. This is the
empty-dictionary engine for proteins: the receiver stores nothing and
synthesises the full state by constraint propagation to convergence, not by
forward simulation. \shk{observe\ P} then forces the measurement now.
\end{definition}

\begin{remark}[The verbs are one verb]
\label{rem:one-verb}
\textsc{Individuate}, \textsc{Contact}, \textsc{Close}/\textsc{Fold},
\textsc{Propagate}, \textsc{Catalyze}, \textsc{Mutate}, and
\textsc{Complete} are the cut at arities one, two, closure, and chain. The
standard library is therefore not seven independent operations but one
operation (the cut) exposed at several arities and readings, exactly as
\Cref{sec:intro} claims and as the underlying theory requires. A reader who
understands the cut understands the whole language.
\end{remark}

% =====================================================================
\section{Core IR and the two-target model}
\label{sec:ir}
% =====================================================================

Shakespeare compiles, from one front end, to a typed core intermediate
representation, and thence to two back ends. This section fixes the IR and
proves the two back ends agree.

\subsection{The core IR}

\begin{definition}[Cut-IR]
\label{def:ir}
\emph{Cut-IR} is a typed sequence of cut instructions over a contact-graph
state carrying a receiver. Each instruction is one of: $\textsc{Ind}(\kappa,a)$
(individuate a leaf), $\textsc{Con}(a,b,g)$ (guarded contact),
$\textsc{Cls}(X,\bar a)$ (close/complex), $\textsc{Fld}(s)$ (fold),
$\textsc{Prp}(x,P,\textit{adm})$ (propagate/track),
$\textsc{Cat}(s,E)$ (catalyze), $\textsc{Mut}(P,m)$ (mutate),
$\textsc{Cmp}(P,\bar o)$ (complete), $\textsc{Obs}(v)$ (force measurement).
Each instruction carries its result type and floor (\Cref{sec:types}) and,
when executed, returns an accountable value and increments the cut count.
Cut-IR is back-end-agnostic: it names cuts, not how they are realised.
\end{definition}

The front end (lexer, parser, type checker of
\Cref{sec:lexical}--\ref{sec:types}) is shared by both targets and lowers a
program to Cut-IR. Only the realisation of a cut instruction differs
between targets.

\subsection{Target R (Rust): cuts as minimum cuts}

\begin{definition}[Rust realisation]
\label{def:rust-target}
In the Rust back end, each Cut-IR instruction is realised by an exact
computation on the contact graph: $\textsc{Ind}$, $\textsc{Con}$,
$\textsc{Cls}$ compute the relevant minimum cut by maximum flow
\citep{fordfulkerson1956,cormen2009}; $\textsc{Fld}$ delegates to the
validated \shk{levinthal-folding} Kuramoto integrator and coupling-spectrum
chain; $\textsc{Prp}$ and $\textsc{Cat}$ enumerate and evaluate admissible
propagations exactly; $\textsc{Cmp}$ iterates constraint propagation to its
fixed point; residues are exact edge-weight sums. This is the reference
semantics: it computes $w(\cut(\cdot))$ to the floating-point precision of
the host.
\end{definition}

\subsection{Target S (TypeScript/WebGL): cuts as render passes}

\begin{definition}[Shakespeare/WebGL realisation]
\label{def:ts-target}
In the TypeScript/WebGL\,2 back end---the online Shakespeare tool---each
Cut-IR instruction is realised as a GPU render pass over the contact-graph
state encoded in textures: a fragment shader evaluates the cut at each
texel and writes the framebuffer, whose contents \emph{are} the accountable
result (rendering--measurement identity, \citealt{sachikonye2025sbs}). The
domain-specific observation hook (\shk{computePartitionValue}) supplies the
protein partition assignment; the remaining pipeline stages (vertex,
interference, quality, display) are universal. Minimum cuts and folding
spectra are computed by GPU-resident kernels; residues are read back from
the framebuffer. This realisation runs in any WebGL\,2 browser in constant
GPU memory.
\end{definition}

\subsection{Target equivalence}

\begin{theorem}[Two-target equivalence up to the floor]
\label{thm:target-equiv}
Let $p$ be a well-typed Shakespeare program with ambient floor $\flo$. Let
$v_R$ and $v_S$ be the values it yields under the Rust and WebGL back ends
respectively. Then $v_R$ and $v_S$ have the same type and the same cut
structure, and their residues agree up to the floor:
\[
\abs{\res(v_R)-\res(v_S)}\;\le\;\flo .
\]
The two back ends never disagree about any distinction Shakespeare can
express, because Shakespeare cannot express a distinction finer than $\flo$
(\Cref{thm:no-zero-value}).
\end{theorem}

\begin{proof}
Both back ends realise the same Cut-IR (\Cref{def:ir}), so they perform the
same sequence of cut instructions with the same types and the same cut
count (the front end is shared). They differ only in how each cut's residue
is computed: Rust by exact max-flow (and the reference folding
integrator), WebGL by GPU render-readback. The GPU realisation computes the
same minimum cut and coupling spectrum to its texel/precision resolution;
its readback error is bounded by the coarsest resolvable boundary, which is
the floor (no Shakespeare value is individuated below $\flo$,
\Cref{thm:no-zero-value}). Hence the residues agree to within $\flo$, and
since Shakespeare admits no sub-floor distinction, the two results are
observationally identical at the language level.
\end{proof}

\begin{remark}[Why two targets, not one]
\label{rem:two-targets}
The split is deliberate and the equivalence theorem is what makes it safe.
The WebGL target is fast, in-browser, and renders cuts as measurements
directly---ideal for the interactive Shakespeare sandbox. The Rust target
is exact and authoritative---ideal as the reference against which the online
tool is validated. \Cref{thm:target-equiv} guarantees the online tool is
never \emph{wrong} relative to the reference; it is at most less precise,
and only below the floor, where the theory forbids meaning anyway.
\end{remark}

% =====================================================================
\section{Worked plays}
\label{sec:examples}
% =====================================================================

\subsection{Fold a protein}

\begin{lstlisting}[caption={Folding crambin (a closure of residue-leaf cuts).}]
-- crambin.shk
receiver bio              -- fix the biomolecular receiver R_bio
floor 1.0e-4              -- nothing is free; ambient floor

C := fold "1CRN"          -- run the folding morphism chain
observe C                 -- force the measurement now
-- C : Fold @ 1.0e-4
--   contact map  N x N, secondary-structure annotated
--   helices 3, sheets 2   (from coupling eigenstructure)
--   native = unique Kuramoto-energy minimum
\end{lstlisting}

\subsection{Run a catalytic cycle}

\begin{lstlisting}[caption={The cytochrome CYP3A4 seven-state cycle; the Miracle Principle licenses Compound I.}]
-- cyp3a4.shk
receiver bio
floor 3.7e-4              -- R_bio floor estimate

-- build the resting-state receiver tree (PDB 1TQN)
E := complex CYP3A4(
       cofactor "heme",
       residue "CYS442",  -- proximal thiolate ligand
       solvent "H2O_axial"
     ) by CYS442

O2 := substrate "O2"

-- turn dioxygen over through the enzyme; admissible iff the cycle
-- S-entropy-converges (substrate -> product, cycle closes).
-- Compound I (state 6) is locally infeasible but licensed by closure.
cycle := track O2 in E
           with reps mass, charge, spin   -- representation switching
           until converge
           yield amalgamation

observe cycle             -- the amalgamation IS the accounted cycle
assert cycle.count == 7  emit "cycle did not close in seven states"
\end{lstlisting}

\subsection{Score a mutation}

\begin{lstlisting}[caption={A point mutation as a re-cut of the coupling graph.}]
-- variant.shk
receiver bio
floor 1.0e-4

wt  := fold "1CRN"
mut := mutate wt by VAL -> GLY at 8   -- re-cut K_ij at residue 8

-- variant effect = change in folding convergence / eigenstructure
observe mut
assert mut.coherence >= wt.coherence  emit "mutation destabilises fold"
\end{lstlisting}

\begin{remark}[The same cut, several jobs]
\label{rem:same-keyword}
Across the three plays, \shk{fold} closes a chain, \shk{complex} closes an
assembly, \shk{track}/\shk{catalyze} propagate a residue chain, and
\shk{mutate} re-cuts---all the same primitive at different arities and
readings (\Cref{rem:one-verb}). \shk{observe} forces the measurement; the
result is the accountable value, never a simulation of it.
\end{remark}

% =====================================================================
\section{Implementation requirements}
\label{sec:impl}
% =====================================================================

Both targets must implement, from this specification:

\begin{enumerate}[leftmargin=1.6em,itemsep=2pt]
\item The shared front end: lexer (\Cref{sec:lexical}), parser to the
grammar (\Cref{def:grammar}), and accountability type checker
(\Cref{sec:types}) enforcing \Cref{thm:no-zero-value}, over the five leaf
classes and the process types.
\item Lowering to Cut-IR (\Cref{def:ir}).
\item The operational semantics of \Cref{sec:opsem}: a contact-graph state
carrying the receiver, a monotone cut count, and the
leaf/contact/close/fold/track/catalyze/mutate/complete reductions, with
evaluation mutating state (measurement, \Cref{cor:measure}).
\item The standard-library bindings of \Cref{sec:stdlib} to the framework
operations \textsc{Individuate}, \textsc{Contact}, \textsc{Close},
\textsc{Fold}, \textsc{Propagate}, \textsc{Catalyze}, \textsc{Mutate},
\textsc{Complete}.
\item A back end: exact max-flow minimum cut plus the \shk{levinthal-folding}
integrator (Rust, \Cref{def:rust-target}), or the GPU render-pass
realisation with the protein observation hook (TypeScript/WebGL,
\Cref{def:ts-target}).
\end{enumerate}

A conformance suite must check \Cref{thm:target-equiv}: every worked play
yields type- and structure-identical results on both targets with residues
agreeing to within the floor. The WebGL target is the online Shakespeare
tool; the Rust target is its oracle.

\begin{remark}[Reference interpreter]
\label{rem:reference}
A direct interpreter of \Cref{sec:opsem} over the exact (max-flow) back end,
reusing the already-validated framework operations as the standard-library
bindings, constitutes the simplest conforming implementation and serves as
the oracle for both production targets.
\end{remark}

% =====================================================================
\section{Conclusion}
\label{sec:conclusion}
% =====================================================================

We have specified \emph{Shakespeare}, a language whose single primitive is
the cut---the act of individuating a part of a bounded whole at the
irreducible cost $\flo>0$---and whose value model is the biomolecular
receiver $\Recv$, so that every value is an S-expression over the five-class
leaf algebra and every evaluation is a run of the receiver's typed morphism
chain. Every operation, from generating a residue leaf to turning a
substrate over through an enzyme, is that primitive at a different arity and
reading: folding is a track over Kuramoto-synchronised residue leaves,
catalysis and measurement are the same cut differing only in whether the
residue's causal link is resolved, and \shk{track\,\ldots\,until\,converge}
is the causal propagation table evaluated over proteins. The
accountability type system makes the floor a static invariant: no
Shakespeare value can claim the forbidden sharp cut. The operational
semantics makes evaluation measurement rather than simulation, with a
monotone cut count as the program's intrinsic clock and a track's
intermediate steps free while only terminal $S$-entropy convergence is
judged---so a catalytic cycle may pass through a transient Compound~I under
the Miracle Principle. The standard library binds each verb to a validated
operation of the framework, so a Shakespeare program is an executable
protein-process paper. And the two-target model---one shared, typed front
end lowering to a back-end-agnostic Cut-IR, realised exactly in Rust and by
GPU rendering in the online WebGL tool---gives both an authoritative
reference and an interactive tool that, by \Cref{thm:target-equiv}, never
disagree about anything the language can express. The protein is a
protocol; a Shakespeare program is its performance, and to run it is to
measure the protein, once, at the floor.

\bibliographystyle{plainnat}
\bibliography{references}

\end{document}
